\section{Abstract}

Pipelining is the predominant microarchitectural technique for increasing processor throughput in RISC designs by overlapping the execution of multiple instructions. Although the literature reports numerous implementations of RV32IM processors on FPGAs, both single-cycle and five-stage pipelined, existing studies typically evaluate only one microarchitecture or rely on heterogeneous platforms and measurement protocols, leaving a gap in directly comparable, reproducible empirical evidence that contrasts both variants of the same core under controlled conditions.
This project addresses that gap by designing and implementing in SystemVerilog a RISC-V RV32IM processor in two microarchitectures, single-cycle and five-stage pipeline with data and control hazard resolution, on the Intel Cyclone V DE1-SoC FPGA. Functional correctness of both implementations is verified using the riscv-tests suite and cocotb-based testbenches covering the full RV32IM instruction set. Through homogeneous synthesis runs with Intel Quartus Prime and a unified measurement protocol, the two designs are experimentally compared across maximum operating frequency, effective CPI, throughput, and utilization of logic elements, memory blocks, and DSP blocks. An additional contribution is a pipeline state visualization module integrated into the same board, designed as a teaching aid for computer architecture courses.
The pipelined microarchitecture is expected to achieve higher throughput than its single-cycle counterpart by sustaining a higher clock frequency, even accounting for the CPI penalty introduced by hazards. The results will provide quantitative, replicable evidence on the real design trade-offs between both microarchitectures on a mid-range FPGA platform widely used in educational settings.

\textbf{Keywords:} RISC-V, RV32IM, pipeline, single-cycle, FPGA, DE1-SoC, SystemVerilog, hazards, functional verification, performance.